\documentclass{article}
\usepackage[utf8]{inputenc}
\usepackage[margin=1in]{geometry}
\usepackage{times}
\usepackage{graphicx}
\usepackage{natbib}
\usepackage{hyperref}

\title{Adaptive Hierarchical Planning with Failure Detection and Strategy Switching}
\author{Mudit Baid \\ Stanford University \\ \texttt{muditbaid@stanford.edu}}
\date{}

\begin{document}
\maketitle

\section{Problem Statement and Motivation}

Large language models struggle with complex, multi-step planning tasks. When a plan fails midway, current systems either restart from scratch or continue down a bad path. They lack the ability to detect failures early and adapt.

This project addresses two gaps. First, LLMs do not decompose problems hierarchically. They generate flat sequences without distinguishing high-level goals from low-level actions. Second, LLMs commit to a single reasoning strategy regardless of subtask type. A math subtask and a commonsense subtask need different approaches, but current systems treat them uniformly.

The goal is to build a planning system that decomposes tasks hierarchically, selects appropriate strategies per subtask, verifies constraints at each step, and recovers from failures through rollback or strategy switching.

\section{Related Work}

HTN planners decompose tasks into subtasks using predefined schemas \citep{erol1994htn}. Recent work applies LLMs to plan generation, but most produce flat action sequences \citep{huang2022language}.

Chain-of-Thought prompting improves reasoning via intermediate steps \citep{wei2022chain}. Tree-of-Thought extends this by exploring multiple paths \citep{yao2023tree}. However, these methods do not adapt strategy based on problem type or recover from dead ends.

Plan verification remains understudied. Some work uses LLMs to check feasibility \citep{valmeekam2023planning}, but integration with rollback is limited. This project combines hierarchical decomposition, adaptive strategy selection, constraint verification, and failure recovery.

\section{Proposed System Design}

The system has five modules:

\textbf{Task Decomposer.} Breaks a complex task into a tree of subgoals. The root is the original task. Children are subtasks that achieve the parent. Leaves are atomic actions.

\textbf{Subtask Classifier.} Labels each subtask by type: arithmetic, logical, commonsense, spatial, or procedural. This determines which reasoning strategy to apply.

\textbf{Strategy Executor.} Runs the appropriate strategy per type: CoT with calculator for arithmetic, ToT for logical, retrieval-augmented prompting for commonsense, state tracking for spatial, precondition checks for procedural.

\textbf{Constraint Verifier.} After each step, checks that output satisfies constraints: logical consistency, resource limits, domain rules. Violations trigger the failure handler.

\textbf{Failure Handler.} When verification fails, decides between: (1) rollback to parent and try alternative decomposition, (2) switch strategy for current subtask, (3) propagate failure upward.

\section{Evaluation Plan}

\textbf{Datasets.} ALFWorld (household tasks), BlocksWorld (classical planning), and a custom benchmark with mixed reasoning types.

\textbf{Baselines.} Flat CoT, Flat ToT, ReAct, and ablations without rollback or strategy switching.

\textbf{Metrics.} Task success rate, steps to completion, rollback frequency and success, strategy switch impact.

\textbf{Analysis.} When does decomposition fail? Which subtask types are hardest? When does rollback help versus hurt?

\section{Timeline and Group}

Weeks 1-2: Decomposer and classifier. Weeks 3-4: Strategy executor. Week 5: Verifier and failure handler. Weeks 6-7: Experiments and ablations. Week 8: Analysis and report.

This project is conducted solo by Mudit Baid.

\bibliographystyle{plainnat}
\begin{thebibliography}{9}
\bibitem[Erol et al., 1994]{erol1994htn}
Erol, K., Hendler, J., and Nau, D. (1994). HTN planning: Complexity and expressivity. \textit{AAAI}.

\bibitem[Huang et al., 2022]{huang2022language}
Huang, W., Abbeel, P., Pathak, D., and Mordatch, I. (2022). Language models as zero-shot planners. \textit{ICML}.

\bibitem[Wei et al., 2022]{wei2022chain}
Wei, J., et al. (2022). Chain-of-thought prompting elicits reasoning in large language models. \textit{NeurIPS}.

\bibitem[Yao et al., 2023]{yao2023tree}
Yao, S., et al. (2023). Tree of thoughts: Deliberate problem solving with large language models. \textit{NeurIPS}.

\bibitem[Valmeekam et al., 2023]{valmeekam2023planning}
Valmeekam, K., et al. (2023). On the planning abilities of large language models. \textit{NeurIPS}.
\end{thebibliography}

\end{document}
